\documentclass[11pt,a4paper]{article}

% =========================================================================
%  RAPPORT — AGENT DE VEILLE CYBERSECURITE (CVE MONITORING)
%  Theme : Veille / SecOps — palette vert-cyan (identité portfolio)
% =========================================================================

\usepackage[utf8]{inputenc}
\usepackage[T1]{fontenc}
\usepackage[french]{babel}
\usepackage{lmodern}
\usepackage[a4paper,margin=2.2cm,top=2.6cm,bottom=2.4cm]{geometry}
\usepackage{xcolor}
\usepackage{fontawesome5}
\usepackage{tikz}
\usetikzlibrary{shapes.geometric,positioning,calc,arrows.meta}
\usepackage{titlesec}
\usepackage[most]{tcolorbox}
\usepackage{listings}
\usepackage{fancyhdr}
\usepackage{booktabs}
\usepackage{tabularx}
\usepackage{array}
\usepackage{colortbl}
\usepackage{enumitem}
\usepackage{microtype}
\usepackage[bookmarks=true,colorlinks=true,linkcolor=accent,urlcolor=accent,citecolor=accent]{hyperref}

% --- Palette Veille (vert / cyan — identité du portfolio) ----------------
\definecolor{primary}{HTML}{0B0F17}
\definecolor{accent}{HTML}{00B37A}
\definecolor{accentdark}{HTML}{007A54}
\definecolor{light}{HTML}{F0FFFA}
\definecolor{midgray}{HTML}{6B7280}
\definecolor{codebg}{HTML}{0D1322}
\definecolor{codefg}{HTML}{E6EEFC}
\definecolor{codekw}{HTML}{4DD2FF}
\definecolor{codestr}{HTML}{9CE6C7}
\definecolor{codecmt}{HTML}{6A8F87}

\titleformat{\section}{\sffamily\Large\bfseries\color{primary}}{\thesection}{0.8em}{}[\vspace{-0.4ex}{\color{accent}\rule{\linewidth}{1.2pt}}]
\titleformat{\subsection}{\sffamily\large\bfseries\color{accent}}{\thesubsection}{0.6em}{}
\titlespacing*{\section}{0pt}{1.6em}{0.8em}

\pagestyle{fancy}
\fancyhf{}
\renewcommand{\headrulewidth}{0pt}
\fancyhead[L]{\footnotesize\sffamily\color{midgray}\faIcon{shield-alt}~Agent de veille CVE — Monitoring automatisé}
\fancyhead[R]{\footnotesize\sffamily\color{midgray}Honorine Kylian}
\fancyfoot[C]{\footnotesize\sffamily\color{midgray}\thepage}

\lstdefinestyle{pystyle}{
  backgroundcolor=\color{codebg},
  basicstyle=\ttfamily\footnotesize\color{codefg},
  keywordstyle=\color{codekw}\bfseries,
  stringstyle=\color{codestr},
  commentstyle=\color{codecmt}\itshape,
  showstringspaces=false, breaklines=true,
  frame=none, framesep=8pt, framerule=0pt,
  xleftmargin=12pt, xrightmargin=12pt,
  columns=fullflexible, inputencoding=utf8, extendedchars=true,
  literate={é}{{\'e}}1 {è}{{\`e}}1 {ê}{{\^e}}1 {à}{{\`a}}1 {ô}{{\^o}}1 {î}{{\^\i}}1 {ç}{{\c{c}}}1 {ù}{{\`u}}1 {â}{{\^a}}1 {É}{{\'E}}1 {È}{{\`E}}1 {Ê}{{\^E}}1 {À}{{\`A}}1 {Ô}{{\^O}}1 {Ç}{{\c{C}}}1 {Ù}{{\`U}}1 {Â}{{\^A}}1
}
\lstset{style=pystyle,language=Python}

\newtcblisting{pybox}[1][]{
  listing only,
  enhanced jigsaw, breakable,
  colback=codebg, colframe=accent,
  arc=2pt, boxrule=0pt, leftrule=3pt,
  left=8pt, right=8pt, top=4pt, bottom=4pt,
  listing options={style=pystyle,language=Python}, #1
}

\newtcolorbox{infobox}[1][]{
  enhanced, colback=light, colframe=accent,
  arc=2pt, boxrule=0pt, leftrule=3pt,
  left=10pt, right=10pt, top=6pt, bottom=6pt,
  fontupper=\small, #1
}
\newtcolorbox{resultbox}[1][]{
  enhanced, colback=accent!10, colframe=accent,
  arc=2pt, boxrule=0pt, leftrule=3pt,
  left=10pt, right=10pt, top=6pt, bottom=6pt,
  fontupper=\small, #1
}

\setlist[itemize]{leftmargin=*,itemsep=2pt,topsep=4pt}
\setlist[enumerate]{leftmargin=*,itemsep=2pt,topsep=4pt}

\begin{document}
\pagenumbering{gobble}

% ---- PAGE DE TITRE ------------------------------------------------------
\begin{tikzpicture}[remember picture, overlay]
  \fill[primary] (current page.north west) rectangle ([yshift=-9cm]current page.north east);
  \fill[accent] ([yshift=-9cm]current page.north west) rectangle ([yshift=-9.25cm]current page.north east);
  % Logo : radar de veille + alerte CVE
  \node[anchor=center] at ([yshift=-4.6cm]current page.north) {%
    \begin{tikzpicture}[scale=1.0]
      \fill[white,rounded corners=4pt] (-2.4,-1.7) rectangle (2.4,1.7);
      \fill[primary,rounded corners=4pt] (-2.4,1.1) rectangle (2.4,1.7);
      \fill[accent] (-2.15,1.4) circle (0.13);
      % Radar concentrique
      \draw[accent,line width=1.1pt] (-0.55,-0.35) circle (1.05);
      \draw[accent,line width=1.1pt] (-0.55,-0.35) circle (0.68);
      \draw[accent,line width=1.1pt] (-0.55,-0.35) circle (0.31);
      \draw[accent,line width=1.3pt] (-0.55,-0.35) -- (0.35,0.55);
      \fill[accent] (-0.55,-0.35) circle (0.05);
      % Points de détection (CVE captées)
      \foreach \x/\y in {-1.1/0.15,-0.05/-0.9,0.15/0.35,-1.0/-0.85} \fill[accentdark] (\x,\y) circle (0.055);
      % Badge alerte
      \node[fill=accent,rounded corners=2pt,inner sep=3pt] at (1.55,0.15) {\faIcon[regular]{bell}~\textcolor{primary}{\scriptsize\bfseries CVE}};
      \node[font=\ttfamily\bfseries\small,white] at (-1.55,1.4) {watchdog};
    \end{tikzpicture}};
  \node[anchor=north, white, font=\sffamily\Huge\bfseries] at ([yshift=-6.7cm]current page.north) {Veille CVE};
  \node[anchor=north, accent, font=\sffamily\Large] at ([yshift=-7.65cm]current page.north) {Agent de monitoring automatisé};
  \node[anchor=north, white!75!primary, font=\sffamily\small] at ([yshift=-8.4cm]current page.north) {Python \textbullet{} Docker Compose \textbullet{} Trivy \textbullet{} Slack};

  \node[anchor=south west, primary, font=\sffamily\small] at ([xshift=2.2cm,yshift=3cm]current page.south west) {\textbf{\textcolor{accent}{Auteur}}};
  \node[anchor=south west, primary, font=\sffamily\small] at ([xshift=2.2cm,yshift=2.5cm]current page.south west) {HONORINE Kylian};
  \node[anchor=south east, primary, font=\sffamily\small] at ([xshift=-2.2cm,yshift=3cm]current page.south east) {\textbf{\textcolor{accent}{Module}}};
  \node[anchor=south east, primary, font=\sffamily\small] at ([xshift=-2.2cm,yshift=2.5cm]current page.south east) {Stage — Infra \& Cybersécurité 2025/2026};
  \fill[accent] ([yshift=1.3cm]current page.south west) rectangle ([yshift=1.4cm]current page.south east);
  \node[anchor=south, midgray, font=\sffamily\footnotesize] at ([yshift=0.6cm]current page.south) {IUT Réseaux \& Télécommunications — Université de La Réunion};
\end{tikzpicture}
\newpage

\pagenumbering{arabic}
\setcounter{page}{1}

{\sffamily
\hypersetup{linkcolor=primary}
\renewcommand{\contentsname}{\textcolor{accent}{\faList~} Sommaire}
\tableofcontents
}
\vspace{0.4cm}

\section{Contexte}

\begin{infobox}[title={\faSatelliteDish~Mission}]
Automatiser la phase \textbf{« identifier »} d'une procédure de maintenance applicative sur l'infrastructure Kubernetes RKE2, inspirée du \textbf{NIST SP 800-40} (gestion des correctifs) et du guide \textbf{ReCyF de l'ANSSI} : construire un agent qui détecte, priorise et notifie les vulnérabilités touchant la stack en production.
\end{infobox}

Sans veille continue et fiable, les trois phases suivantes de la procédure — évaluer, corriger, vérifier — restent aveugles. Ce projet a été mené \textbf{en autonomie complète pendant un stage}, sans accompagnement technique de l'entreprise sur ce sujet précis : identification des sources, conception de l'architecture, logique métier, audit de sécurité de l'outil lui-même.

\medskip
\noindent\textbf{Stack technique :}
\begin{itemize}
  \item \textbf{Python} (\raise0pt\hbox{\textasciitilde}1100 lignes) \textemdash{} logique de récupération, matching et notification
  \item \textbf{Docker Compose} \textemdash{} conteneurisation et déploiement en production
  \item \textbf{Trivy} \textemdash{} scan de vulnérabilités des images Docker
  \item \textbf{Prometheus / Grafana} \textemdash{} observabilité de l'agent
  \item \textbf{Claude (IA générative)} \textemdash{} assistant de développement, sous ma supervision
\end{itemize}

% =========================================================================
\section{Sources de veille}

Cinq sources officielles, interrogées via leurs \textbf{API} plutôt que par scraping — plus rapide, plus stable, respectueux des CGU et indépendant du design des pages.

\renewcommand{\arraystretch}{1.3}
\begin{tabularx}{\linewidth}{@{}l X l@{}}
\toprule
\rowcolor{light} Source & Rôle & Fréquence \\ \midrule
\textbf{CISA KEV} & CVE confirmées activement exploitées & 15 min \\
\textbf{NVD} (NIST) & Base mondiale des CVE, score CVSS & 30 min \\
\textbf{CERT-FR} (ANSSI) & Avis de sécurité France / Europe & 30 min \\
\textbf{GHSA} (GitHub) & Vulnérabilités des dépendances open source & 1 h \\
\textbf{endoflife.date} & Fins de support produits & 1 / jour \\
\textbf{Trivy} (images Docker) & Scan CVE de 9 images de la stack monitoring & 1 / jour \\
\bottomrule
\end{tabularx}

\textbf{Pourquoi des fréquences différenciées :} KEV est la source la plus actionable (exploitation confirmée) donc pollée le plus souvent ; Trivy est volontairement limité à une fois par jour pour ne pas saturer le canal Slack.

% =========================================================================
\section{Architecture \& boucle de veille}

L'agent tourne en \textbf{mode daemon continu} dans un conteneur Docker, chaque source étant pollée selon sa propre fréquence.

\begin{pybox}
SOURCES = {
    "cisa_kev":      {"interval": 15 * 60},
    "nvd":           {"interval": 30 * 60},
    "cert_fr":       {"interval": 30 * 60},
    "ghsa":          {"interval": 60 * 60},
    "endoflife":     {"interval": 24 * 3600},
    "trivy_images":  {"interval": 24 * 3600},
}

def run_forever():
    while True:
        for name, cfg in SOURCES.items():
            if due(name, cfg["interval"]):
                items = fetch(name)          # appel API dedie
                candidates = match_stack(items)   # filtrage TEEO
                for cve in candidates:
                    priority = score_priority(cve)
                    if should_notify(cve, priority):
                        notify_slack(cve, priority)
        sleep(POLL_TICK)
\end{pybox}

\begin{infobox}[title={\faRobot~Rôle de l'IA générative}]
Le code Python a été généré avec l'assistance de \textbf{Claude}. Mon rôle : architecte et testeur — spécifications fonctionnelles, logique métier, validation de chaque itération, correction des bugs découverts en production. L'IA a accéléré l'écriture, la maîtrise du fonctionnement reste entièrement mienne.
\end{infobox}

% =========================================================================
\section{Filtrage \& priorisation}

\subsection{Matching sur la stack surveillée}

Environ 30\,000 CVE sont publiées chaque année \textemdash{} l'immense majorité sans rapport avec l'infrastructure réelle. Le matching combine deux mécanismes :

\begin{itemize}
  \item \textbf{Matching CPE} (Common Platform Enumeration) \textemdash{} identifiants normalisés NVD, fiable et sans faux positif ;
  \item \textbf{Matching textuel} \textemdash{} recherche des \raise0pt\hbox{\textasciitilde}179 composants de la stack (Kong, Keycloak, Vault, PostgreSQL, Grafana, Kubernetes\ldots) dans titre/description, utilisé en repli quand le CPE est absent.
\end{itemize}

\subsection{Grille de priorité P0--P4}

Inspirée de la directive \textbf{CISA BOD 22-01}, la priorité combine score CVSS, statut KEV et exposition du composant (public via Kong/Keycloak, ou interne au cluster).

\renewcommand{\arraystretch}{1.3}
\begin{tabularx}{\linewidth}{@{}>{\bfseries}p{1cm} X l@{}}
\toprule
\rowcolor{light} Prio. & Critère dominant & SLA \\ \midrule
P0 & CVSS élevé + KEV + composant exposé & 24--72 h \\
P1 & CVSS élevé ou KEV, exposition partielle & 14 jours \\
P2 & CVSS moyen, composant interne critique & 30 jours \\
P3 & CVSS moyen, faible exposition & 60 jours \\
P4 & CVSS faible / composant peu critique & 90 jours \\
\bottomrule
\end{tabularx}

% =========================================================================
\section{Déduplication \& anti-flood}

Une même CVE apparaît souvent dans plusieurs sources à quelques jours d'écart. Sans garde-fou, l'équipe recevrait plusieurs alertes redondantes pour la même faille \textemdash{} fatigue d'alerte garantie.

\begin{itemize}
  \item \textbf{Base SQLite} des CVE déjà notifiées \textemdash{} une CVE revue dans une autre source est ignorée par défaut ;
  \item \textbf{Re-notification sur upgrade} \textemdash{} si le statut passe en KEV ou si la priorité augmente, une nouvelle alerte part avec un bandeau \og PRIORITÉ ÉLEVÉE \fg{} ; l'inverse (priorité qui baisse) ne re-notifie jamais ;
  \item \textbf{Mode baseline} au premier démarrage \textemdash{} ingestion silencieuse des 500 dernières CVE, sans notification, pour éviter un flood initial de 50+ messages (pattern standard : warmup Prometheus, learning mode SIEM/EDR).
\end{itemize}

% =========================================================================
\section{Sécurité de l'agent \& observabilité}

\begin{resultbox}[title={\faLock~Un outil de sécurité doit être sécurisé lui-même}]
Protection SSRF sur les liens suivis, exécution en utilisateur \textbf{non-root}, validation stricte des noms d'image Docker (anti-injection via Trivy), fail-fast au démarrage si une variable d'environnement critique manque.
\end{resultbox}

L'agent expose \textbf{14 métriques Prometheus} couvrant les quatre Golden Signals SRE (latence, trafic, erreurs, saturation), visualisées dans un dashboard Grafana à 7 sections. Une CLI (\texttt{docker exec}) permet diagnostic et tests : \texttt{-{}-stats}, \texttt{-{}-history}, \texttt{-{}-test-pipeline}.

% =========================================================================
\section{Résultats concrets}

Après deux mois de production continue (24 h/24) :

\begin{itemize}
  \item Remontée d'une CVE \textbf{Palo Alto} en priorité P0 via le flux KEV ;
  \item Détection de la faille \textbf{ArgoCD ServerSideDiff} (CVSS 9.6), ayant déclenché une planification de mise à jour ;
  \item Détection d'images Docker cumulant plus de 100 CVE HIGH, motivant un changement de stratégie d'image de base ;
  \item Auto-surveillance : alerte Slack si l'agent lui-même reste silencieux plus de 4 h.
\end{itemize}

\section{Synthèse technique}

\renewcommand{\arraystretch}{1.3}
\begin{tabularx}{\linewidth}{@{}>{\bfseries}p{2.6cm} X@{}}
\toprule
\rowcolor{light} Bloc & Rôle \\ \midrule
Fetchers API      & Interrogation des 5 sources CVE à fréquence dédiée \\
Trivy runner      & Scan quotidien de 9 images Docker via subprocess \\
Matching engine   & Filtrage CPE + textuel sur la stack surveillée \\
Priority engine   & Score P0--P4 (CVSS + KEV + exposition) \\
Dedup store       & SQLite, re-notification sur upgrade uniquement \\
Notifier          & Webhook Slack + backup SMTP \\
Observability     & 14 métriques Prometheus, dashboard Grafana, CLI \\
\bottomrule
\end{tabularx}

\section{Limites \& évolutions}

Liste des 9 images Trivy maintenue manuellement \textemdash{} l'évolution naturelle serait un \textbf{Trivy Operator} déployé dans le cluster pour une couverture exhaustive automatique. Absence de ticketing (Jira/Asana) pour le suivi du cycle de vie des vulnérabilités, attendu par \textbf{NIS2}. Pas de détection runtime (\textbf{Falco}) ni d'audit de configuration Kubernetes (\textbf{Kubescape}) \textemdash{} deux angles morts identifiés pour une industrialisation future.

\section{Conclusion}

Ce projet démontre qu'un outil de veille cybersécurité robuste peut être conçu et livré en autonomie, en s'appuyant sur l'IA générative comme accélérateur de développement sans jamais en déléguer la logique métier ni la responsabilité. Le résultat : un agent qui transforme un flux brut de \raise0pt\hbox{\textasciitilde}30\,000 CVE annuelles en une poignée d'alertes réellement actionables.

\vspace{0.5cm}
\begin{center}
{\sffamily\itshape\color{midgray}\large
\og Sans veille continue, la correction arrive toujours trop tard. \fg{}}
\end{center}

\end{document}
